\section{Other dynamics models}\label{sec:other_dynamics}
These models modify the equations of motion of \autoref{sec:rocket2d}, or extend them by a degree
of freedom, rather than supply a term of them. All are off by default.

\subsection{Launch rail}\label{sec:rail}
On the launch rail the rocket can only slide along it: the rail buttons react every force component
transverse to the rail and prevent rotation. This is a \emph{kinematic constraint}, not a new
force: the assembled acceleration is projected onto the rail, and the rotation is removed.

The rail direction is fixed and is the direction the nose initially points. With $\theta_{rail}$
the rail elevation, an input set consistently with the initial $\theta_r$,
\begin{equation} \label{eq:rail_direction}
    \hat{e}_{rail} = \cos\theta_{rail}\,\hat{x}_0 - \sin\theta_{rail}\,\hat{z}_0,
\end{equation}
that is, $\hat{x}_r$ at $\theta_r=\theta_{rail}$, frozen. The distance travelled from a launch
point at the origin is the projection of the position onto it,
\begin{equation} \label{eq:rail_distance}
    s_{rail}(\vec{x}_r) = \vec{x}_r\cdot\hat{e}_{rail}
        = x\cos\theta_{rail}-z\sin\theta_{rail}.
\end{equation}
The rocket is on the rail while $s_{rail}<L_{rail}$ (a state predicate), and there the equations
of motion are replaced by
\begin{gather}
    \vec{a}_r = \left[\left(\frac{\sum\vec{F}}{m}
        + \vec{a}_{Coriolis}^{\,2D}\right)\cdot\hat{e}_{rail}\right]\hat{e}_{rail}
        \label{eq:rail_constraint} \\
    \omega_r = 0, \qquad \varepsilon_r = 0 \label{eq:rail_attitude}
\end{gather}
with $\sum\vec{F}$ the sum of \autoref{eq:eom_trans} and $\vec{a}_{Coriolis}^{\,2D}$ the term of
\autoref{eq:coriolis_2d} if enabled. The component of $\vec{a}_r$ normal to the rail is identically
zero, and since $\vec{x}_r(0)$ lies on the rail line with $\vec{v}_r(0)=\vec{0}$, position and
velocity stay on it with no correction step. In the integrator, \autoref{eq:rail_attitude} reads
$\mathrm{d}\theta_r/\mathrm{d}t=\mathrm{d}\omega_r/\mathrm{d}t=0$, so $\theta_r$ and $\omega_r$ keep
their initial values. The pitching moment is reacted by the rail rather than allowed to act.
\figtodo{Diagram (TikZ): rail direction $\hat{e}_{rail}$, elevation $\theta_{rail}$, length
$L_{rail}$ and the distance $s_{rail}$ along it.}

\paragraph{Exit and reverse motion.} $L_{rail}=0$ is the default and the off switch. Past rail
exit ($s_{rail}\geq L_{rail}$) the free-flight equations apply again, with no transition term.
Nothing holds the rocket on the pad before thrust exceeds weight.
If the thrust along the rail is below the weight component, \autoref{eq:rail_constraint} gives a
negative $\vec{a}_r\cdot\hat{e}_{rail}$ and the rocket slides \emph{backwards} down the rail
($s_{rail}<0$) until thrust catches up. A real pad reacts weight until lift-off, and RASAero II
corrected exactly this case \cite{rasaerorel}. The fix is to clamp $\vec{a}_r\cdot\hat{e}_{rail}\geq0$
while $s_{rail}\leq0$ and $\vec{v}_r\cdot\hat{e}_{rail}\leq0$.
\restodo{decision: adopt the clamp; it would change existing results for low thrust-to-weight
starts}
\comptodo{no rail ($L_{rail}=0$) against a rail of the demo launch site on the demo rocket; the
table gives rail-exit speed, apogee and downrange. Generated by the simulator.}
The key output is the rail-exit speed $|\vec{v}_r|$, which decides whether wind can tip the
airframe before it can weathercock. There is no rail friction, flex or tip-off: the transverse and
rotational freedoms are restored instantly at $s_{rail}=L_{rail}$.

\subsection{Recovery}\label{sec:recovery}
Once a canopy is deployed the attitude is frozen exactly as on the rail, $\omega_r=0$,
$\varepsilon_r=0$, and only the translational equation changes. Deployment is instantaneous and
total, and the canopy is modelled by its drag alone:
\begin{equation} \label{eq:chute_eom}
    m\vec{a}_r = \vec{W} + \vec{T} + \vec{A}_f + \vec{R}_{ch} + m\,\vec{a}_{Coriolis}^{\,2D},
\end{equation}
with $\vec{R}_{ch}$ the canopy drag below and $\vec{a}_{Coriolis}^{\,2D}$ the term of
\autoref{eq:coriolis_2d}, present only when the Earth-rotation correction is enabled. The
orientation is held at the pitch angle at deployment, and the pitching moment is not integrated.
Freezing is chosen over an artificial damping term, which would need a time constant no source
fixes.
\simtodo{the Coriolis term is dropped in the code under canopy, although it is included on the rail
and in free flight}
\restodo{decision, no literature: in the code only $\mathrm{d}\theta_r/\mathrm{d}t$ and
$\mathrm{d}\omega_r/\mathrm{d}t$ are set to zero, so the stored $\omega_r$ keeps its non-zero value
at deployment and the reported pitch rate is meaningless under canopy. Reset it to zero?}

\paragraph{Deployment.} A function of the state, not of a detected event: either a timer,
$t\geq t_{dep}$, or apogee, taken as $v_z>0$ (descending, $z$ being positive down). The latter is
reconstructed from the state because the right-hand side cannot consult an event not yet
finalised, so a pre-apogee transient in $v_z$ --- possible on an unstable trajectory --- would arm
deployment early. The apogee trigger is ignored for the first $0.1\,\mathrm{s}$, which covers a
tail-first sag off the pad when the initial thrust is below the weight; later transients are not
guarded against. These are the only two triggers; the predicate is not latched, but both are
monotone on a nominal flight. While the rocket is on the rail the rail mode takes precedence.
\simtodo{no altitude (or velocity/dynamic-pressure) deployment trigger exists; a
main-chute-at-altitude trigger and a drogue-plus-main sequence are not implemented}

\paragraph{Force.} The canopy drag is computed directly in $\mathcal{F}_0$, like weight and unlike
the body-axis aerodynamic force:
\begin{equation} \label{eq:chute_force}
    \vec{R}_{ch} = -\tfrac12\rho\,C_{D,ch}S_{ch}\,v_{rel}\,\vec{v}_{rel},
\end{equation}
with $S_{ch}$ the canopy reference area and $C_{D,ch}$ its drag coefficient. Three properties are
the whole content of the model:
\begin{itemize}
    \item It is \emph{natively inertial}, not routed through the body axial/normal decomposition. A
    round canopy is axisymmetric and its drag opposes the relative-wind \emph{vector}, not the nose
    direction, which under canopy is frozen at whatever attitude deployment left. Routing it through
    $\hat{x}_r$ would make the descent rate depend on that attitude. $C_{D,ch}$ is therefore a drag
    coefficient on $\hat{v}_{rel}$, not the axial $C_x$, with its own reference area $S_{ch}$.
    \item It is \emph{added to} the airframe's forces: weight, any residual thrust and the body-axis
    aerodynamic force all continue to act, evaluated at the frozen $\theta_r$.
    \item It is written as $v_{rel}\,\vec{v}_{rel}$ rather than $\tfrac12\rho v_{rel}^2C_{D,ch}S_{ch}\,
    \hat{v}_{rel}$, so it has no $0/0$ direction at $v_{rel}=0$.
\end{itemize}
A run can be checked against the terminal descent speed, where $\vec{R}_{ch}$ balances weight,
approached rather than attained since $\rho$ varies:
\begin{equation} \label{eq:chute_terminal}
    v_{term} = \sqrt{\frac{2mg}{\rho\,C_{D,ch}S_{ch}}}.
\end{equation}
$S_{ch}=0$ is the default, as is ``no recovery model''; either alone leaves this subsection inert.
Inflation is instantaneous and total, with no opening transient, shock or reefing, and the frozen
attitude replaces the pendulum swing under the risers.
\figtodo{Plot (simulator output, a named demo scenario): altitude and vertical speed against time
through deployment, with the terminal speed $v_{term}$ of \autoref{eq:chute_terminal} as a
reference line.}

\paragraph{Lumped alternative.} Rocket and canopy can instead be one body with a combined drag
area, $\vec{R}_{ch}=-\tfrac12\rho\,(C_DS)_{sys}\,v_{rel}\,\vec{v}_{rel}$ with the airframe's own
force dropped, where $(C_DS)_{sys}$ is a constant or a table against speed. For a steady descent
this is sound and is the usual descent-rate and drift estimate: the drag areas of canopy and
vehicle add, $v_{term}=\sqrt{2mg/\sum C_DS}$ \cite{whitmore}, whereas OpenRocket and RocketPy take
the drag from the recovery device alone and ignore the airframe \cite{ortech,hprrec}. It cannot
represent the inflation transient and opening shock, the pendulum motion, added mass or the riser
dynamics, which need a drag area growing in time and a second body.

Neglected, with the place to look them up in Knacke \cite{knacke}: the opening transient (drag area
growing from zero over an inflation time, ch.\ 5.4.3--5.4.4), the opening-shock load on the airframe
and shock cord (5.4.6), a reefed (5.6) or two-stage drogue and main sequence (ch.\ 7), the pendulum
and coning motion of the body under the risers (5.3), and the added-mass effect of the air carried
with the canopy. Knacke's rules come from large canopies ($Re\sim10^7$) and under-predict the
opening force and the inflation time of small ones \cite{jatm}.
\restodo{added-mass effect of the canopy: no source found}
\comptodo{the two deployment triggers (timer and apogee) on the demo rocket with a parachute; the
table gives deployment altitude and speed, time of descent, landing distance and peak deceleration.
Generated by the simulator.}

\subsection{Air brake}\label{sec:airbrake}
A switchable extra drag, rigidly mounted to the body, used to modulate drag during the unpowered
coast. Its aerodynamics is part of the rocket's own: the deployment fraction $f_b\in[0,1]$ is one
more input to the coefficients of \autoref{sec:aero}, in one of two forms.
\begin{itemize}
    \item \textbf{Own component.} The brake is a component $b$ of the multi-component build-up
    (\autoref{eq:aero_multi}) with its own reference area $S_b$ and length $c_b$, coefficients
    $C_{x,b},C_{z,b},C_{m,b}$ and reference point $pp_b$, its force scaled by the exposed fraction,
    $S(f_b)=f_bS_b$:
    \begin{equation} \label{eq:airbrake_component}
        \vec{A}_{f,b} = -f_b\,qS_b\left(C_{x,b}\hat{x}_r + C_{z,b}\hat{z}_r\right), \qquad
        \vec{A}_{m,b} = -f_b\,qS_bc_b\,C_{m,b}\hat{y}_r.
    \end{equation}
    It acts at $pp_b$, so \autoref{eq:aero_multi} also gives the pitching moment of its arm.
    \item \textbf{Whole-rocket tables with a deployment axis.} The coefficient tables take $f_b$
    as an extra axis, $(C_x,C_z,C_m,l_p)=\mathrm{interp}(\alpha,M,f_b;\ \text{table})$, from CFD or
    wind-tunnel data with bilinear interpolation in $M$ and $f_b$; this is common practice
    \cite{farha}. A brake can shift the centre of pressure through blanking of the fin flow (one
    design's static margin fell from $3.37$ to $0.25$ calibers at $M=0.2$ \cite{farha}), so $C_z$
    and $C_m$ should be tabulated against $f_b$ too, not only $C_x$.
\end{itemize}
A flat plate with a constant drag coefficient is the usual simplification of the first form.
The simplest case of the first form has only axial drag, $C_{z,b}=C_{m,b}=0$ and $pp_b=p$, with
$C_{x,b}=C_{D,b}$: one additive axial term
\begin{equation} \label{eq:airbrake_force}
    F_b = q\,C_{D,b}\,S_b\,f_b, \qquad \Delta C_x = \frac{C_{D,b}S_b}{S}\,f_b,
\end{equation}
multiplied like $C_x$ by the axial-flow factor \autoref{eq:axial_factor}, so that it opposes the
flow tail-first. $S_b=0$ is the default.
\simtodo{only this additive axial term is implemented; neither the own-component form nor the
$f_b$ table axis is}
\restodo{no normal-force ($C_N$) data for deployed brakes found in the literature}

Two orthogonal choices set $f_b$: a \emph{law}, which decides how far open the brake is, and a
\emph{trigger}, which decides from when it may open.

\subsubsection{Deployment law}\label{sec:airbrake_law}
Three interchangeable variants:
\begin{itemize}
    \item \textbf{None} (default): $f_b\equiv0$.
    \item \textbf{Proportional}: P control on \emph{predicted} apogee,
    \begin{equation} \label{eq:airbrake_p}
        f_b = \mathrm{clip}\!\left(k_p\left(\hat{h}_a-h_{tgt}\right),\,0,\,1\right), \qquad
        \hat{h}_a = h + \frac{v_z^2}{2g},
    \end{equation}
    with $h_{tgt}$ the commanded apogee and $k_p$ the gain. $\hat{h}_a$ is the altitude the rocket
    would reach if drag vanished now, deliberately an \emph{upper} bound that ignores all further
    drag, including the brake's own. A positive error (predicted overshoot) commands more brake.
    The law is gated to the unpowered-ascent window, returning $0$ whenever $T(t)>0$ or $v_z\geq0$;
    that window belongs to the law and is not the trigger below.
    \item \textbf{Full}: $f_b=1$ from the trigger instant onward and $0$ before: a constant area,
    opened once. A sizing study solves against this variant: with $f_b$ a step, a target apogee
    fixes either $C_{D,b}S_b$ or the deployment time, one unknown against one equation.
\end{itemize}
Every law is \emph{memoryless}, re-evaluated from the current state at every right-hand-side call,
since the solver evaluates at arbitrary, repeated and rejected times. A law with memory needs the
sampled mode (\autoref{sec:sampled}); none is written.
\comptodo{the three laws (none, proportional, full) on the demo rocket with an air brake; the
table gives apogee, miss against $h_{tgt}$ and time of maximum $f_b$, the plot the predicted apogee
$\hat{h}_a$, $f_b$ and altitude against time under the proportional law, with $h_{tgt}$ marked.
Generated by the simulator.}

\subsubsection{Deployment trigger}\label{sec:airbrake_trigger}
Which condition arms the brake, chosen independently of the law. Each variant is an \emph{onset}
latch: once true it stays armed.
\begin{itemize}
    \item \textbf{After burnout} (default): $T(t)\leq0$, true from burnout onward.
    \item \textbf{After time}: $t\geq t_{brake}$, with $t_{brake}=0$ (default) meaning deployed from
    $t=0$.
    \item \textbf{After apogee}: $v_z>0$, the canopy's apogee predicate, with the same caveat: a
    pre-apogee transient in $v_z$ would arm it early.
\end{itemize}
Only the \textbf{full} law reads the trigger; \textbf{none} ignores it and \textbf{proportional}
uses its own window. No condition retracts an armed brake, and there is no deployment-rate limit
between $f_b=0$ and $1$; a real actuator has both, and both are future work.

\begin{aitext}
\subsection{Roll and the reaction wheel}\label{sec:roll}
The rocket's own spin about its nose axis $\hat{x}_r$, and an internal flywheel that drives that
spin by exchanging angular momentum with the airframe. This is the one model of this section that
does not modify \autoref{eq:eom_trans} or \autoref{eq:eom_rot} at all: it \emph{adds} a degree of
freedom which, in the pitch plane, is exactly decoupled from both. The roll angle is $\gamma_r$
about $+\hat{x}_r$ by the right-hand rule, the roll rate $\omega_x=\dot\gamma_r$ and the roll
acceleration $\varepsilon_x=\dot\omega_x$, named after the pitch triple $\theta_r,\omega_r,
\varepsilon_r$ of \autoref{sec:frames}. The vehicle's axial moment of inertia about $\hat{x}_r$ is
$I_x$, taken with the wheel \emph{included and locked}.

\subsubsection{Why roll decouples exactly}\label{sec:roll_decoupling}
For an axisymmetric body, $I_{yy}=I_{zz}=I$ with $I$ the pitch inertia of \autoref{sec:mass}, and
Euler's equation for the pitch channel reads
\begin{equation} \label{eq:roll_gyro_pitch}
    I\,\varepsilon_r = \left[\vec{gt}\times\vec{T} + \vec{gp}\times\vec{A}_f
        + \vec{A}_m\right]\cdot\hat{y}_r - \left(I_x-I\right)\omega_z\,\omega_x,
\end{equation}
with $\omega_z$ the yaw rate about $\hat{z}_r$. Planar flight is $\omega_z\equiv0$
(\autoref{sec:assumptions}), so the gyroscopic term vanishes \emph{identically} --- not to first
order, and not for small $\omega_x$, but for any roll rate whatever. Roll cannot drive pitch
without a yaw rate to pair with. A wheel carrying axial momentum $h_w\hat{x}_r$ on a body pitching
at $\omega_r\hat{y}_r$ likewise torques the airframe by
\begin{equation} \label{eq:roll_gyro_wheel}
    \vec{\tau}_{gyro} = -\,\vec{\omega}\times\vec{h}
        = -\left(\omega_r\hat{y}_r\right)\times\left(h_w\hat{x}_r\right)
        = +\,\omega_r h_w\,\hat{z}_r,
\end{equation}
which is a \emph{yaw} torque: no $\hat{y}_r$ component, so it cannot enter
\autoref{eq:roll_gyro_pitch}, and no $\hat{x}_r$ component, so it cannot enter the roll equation
either. Finally, $C_x,C_z,C_m$ and $C_{m,q}$ of \autoref{sec:aero} are functions of $(\alpha,M)$
alone and take no roll angle. So the three channels --- translation, pitch, roll --- share only the
flow state, and the roll equations below can be integrated alongside the others with no term
crossing between them.

What the decoupling \emph{hides} is the yaw that both
\autoref{eq:roll_gyro_pitch} and \autoref{eq:roll_gyro_wheel} produce, which a 6DOF model would
carry and this one cannot represent at all --- exactly as the out-of-plane wind component is
computed and discarded. Two further effects are dropped with it: for a \emph{finned} airframe
$C_z$ and $C_m$ do depend weakly on $\gamma_r$ (a four-fin body at incidence varies with a
$90^\circ$ period, through fin--body interference and the fins' effective dihedral relative to the
$\alpha$ plane), and at high spin there is roll lock-in with the pitch natural frequency and a
Magnus side force.
\restodo{roll-orientation dependence of $C_z,C_m$ for a four-fin airframe: needs a
$C_z(\alpha,M,\gamma_r)$ table axis, for which no data source is identified}

\subsubsection{Equations of motion}\label{sec:roll_eom}
The vehicle is two bodies sharing one axis: the airframe, and a balanced wheel spinning about
$\hat{x}_r$ with its own axial inertia $I_w$ and absolute rate $\omega_w$. Only two torques act
about that axis --- the external aerodynamic moment $L_x$ on the airframe, and the motor torque
$\tau_w$, which by Newton's third law is $+\tau_w$ on the wheel and $-\tau_w$ on the airframe:
\begin{gather}
    \dot\gamma_r = \omega_x \label{eq:roll_kinematics} \\
    \left(I_x-I_w\right)\varepsilon_x = L_x - \tau_w \label{eq:roll_airframe} \\
    I_w\,\dot\omega_w = \tau_w \label{eq:roll_wheel}
\end{gather}
This is the form to integrate, and it is the reason the wheel's \emph{absolute} rate is carried
rather than its relative one: \autoref{eq:roll_wheel} is then a standalone scalar equation with no
$\varepsilon_x$ on its right-hand side. The rate a motor encoder reads is the relative
$\omega_{w/r}=\omega_w-\omega_x$, and the wheel's stored momentum is $h_w=I_w\,\omega_{w/r}$.

Summing \autoref{eq:roll_airframe} and \autoref{eq:roll_wheel} gives the same content as a
momentum balance, which is how the device is best read:
\begin{equation} \label{eq:roll_momentum}
    h_x = I_x\,\omega_x + h_w, \qquad \dot h_x = L_x.
\end{equation}
With no external moment and the vehicle starting at rest, $h_x\equiv0$ and
\begin{equation} \label{eq:roll_exchange}
    \omega_x = -\,\frac{h_w}{I_x}:
\end{equation}
spin the wheel up to momentum $h_w$ and the airframe counter-rotates, using no external torque at
all. That is the whole purpose of a reaction wheel, and \autoref{eq:roll_exchange} is the check a
run is verified against.

\subsubsection{Authority and saturation}\label{sec:roll_saturation}
The device is fictitious unless both of its limits are carried:
\begin{equation} \label{eq:roll_limits}
    \left|\tau_w\right| \leq \tau_{max}, \qquad
    \left|\omega_{w/r}\right| \leq \omega_{w,max}.
\end{equation}
The first is the motor's stall torque. The second is the one that matters: on reaching it the wheel
\emph{saturates}, \autoref{eq:roll_exchange} caps the roll rate it can hold at
$I_w\omega_{w,max}/I_x$, and all authority in that direction is gone until an \emph{external}
moment desaturates it. The speed limit is one-sided in the model --- at saturation the wheel may
still be braked back, only accelerated further is forbidden --- since a two-sided clamp would lock
it at its limit for the rest of the flight. Both limits make $\tau_w$ discontinuous in the state,
exactly as the air brake's $\mathrm{clip}$ does in \autoref{eq:airbrake_p}, and no event is located
at the boundary.
\restodo{$I_x$ of the airframe, and $I_w$, $\tau_{max}$, $\omega_{w,max}$ of a wheel sized for it:
no component is chosen, so no numbers exist}

\subsubsection{Aerodynamic roll moment}\label{sec:roll_aero}
Roll is damped, and can be driven, by the fins. In the convention of \autoref{sec:conventions},
with the non-dimensional roll rate formed like the pitch rate $\hat\omega$ of
\autoref{eq:pitch_damping_moment},
\begin{equation} \label{eq:roll_aero_moment}
    \vec{A}_{m,x} = -qSc\left(C_l + C_{l,p}\,\hat\omega_x\right)\hat{x}_r, \qquad
    \hat\omega_x = \frac{\omega_x c}{2v_{rel}},
\end{equation}
so that $L_x=\vec{A}_{m,x}\cdot\hat{x}_r$, with $C_l$ the static roll-moment coefficient of a
canted fin set and $C_{l,p}>0$ damping, as $C_{m,q}>0$ damps in pitch. The damping term is
evaluated as $\tfrac14\rho v_{rel}\omega_x c^2 C_{l,p}$, algebraically identical but finite at
$v_{rel}=0$, which the literal quotient is not. Both coefficients are zero by
default, which makes $L_x\equiv0$ and leaves \autoref{eq:roll_exchange} --- pure momentum exchange
--- as the model's default behaviour. A motor also produces a roll torque about $\hat{x}_r$ from its
internal vortex flow and nozzle ablation (\autoref{sec:noise}); it belongs in $L_x$, and no
prediction method for it is known.
\restodo{$C_l$ and $C_{l,p}$ for this airframe: the fin cant is not chosen and no roll data exists;
Barrowman-style build-up or CFD would supply both}

\subsubsection{Control law and the constrained phases}\label{sec:roll_law}
The commanded torque is a memoryless function of the state, re-evaluated at every right-hand-side
call for the reason given in \autoref{sec:airbrake_law}. Three variants:
\begin{itemize}
    \item \textbf{None} (default): $\tau_w\equiv0$. The wheel is present but never commanded, so
    the roll that results is whatever the initial condition and $L_x$ produce.
    \item \textbf{Rate null}: $\tau_w=k_d\,\omega_x$, which by \autoref{eq:roll_airframe} opposes
    the roll rate --- the despin law.
    \item \textbf{Angle hold}: $\tau_w=k_p\left(\gamma_r-\gamma_{cmd}\right)+k_d\,\omega_x$, a PD
    hold on a commanded roll angle. There is no integral term: one needs the sampled mode
    (\autoref{sec:sampled}) to carry its accumulator, and none is written.
\end{itemize}
Both gains are positive, and both are clipped by \autoref{eq:roll_limits}.

On the launch rail (\autoref{sec:rail}) the buttons running in the rail prevent roll as they
prevent pitch, so $\omega_x=0$ and $\varepsilon_x=0$ there, alongside
\autoref{eq:rail_attitude}; under a canopy (\autoref{sec:recovery}) roll is frozen the same way,
inheriting that section's choice rather than deserving it --- a real rocket under a canopy spins
freely. In both phases \autoref{eq:roll_wheel} keeps running: the wheel is internal, its reaction
torque is taken by the rail buttons or the risers instead of by the airframe's roll freedom, and a
wheel spun up on the rail therefore reaches rail exit with real stored momentum.
\simtodo{the sampled mode carries no roll channel, so a roll law with memory (integral term,
rate limiting on $\tau_w$, or a measured rather than true $\omega_x$) is not expressible}
\figtodo{Plot (simulator output, a named demo scenario): roll angle $\gamma_r$, roll rate
$\omega_x$ and relative wheel rate $\omega_{w/r}$ against time under the angle-hold law, with the
saturation level $\omega_{w,max}$ of \autoref{eq:roll_limits} as a reference line.}
\comptodo{the three laws (none, rate null, angle hold) on the demo rocket given an initial roll
rate; the table gives the roll angle at apogee, the peak relative wheel rate against
$\omega_{w,max}$, and whether the wheel saturated. Generated by the simulator.}
\end{aitext}
